% Options for packages loaded elsewhere
\PassOptionsToPackage{unicode}{hyperref}
\PassOptionsToPackage{hyphens}{url}
%
\documentclass[
]{article}
\usepackage{amsmath,amssymb}
\usepackage{iftex}
\ifPDFTeX
  \usepackage[T1]{fontenc}
  \usepackage[utf8]{inputenc}
  \usepackage{textcomp} % provide euro and other symbols
\else % if luatex or xetex
  \usepackage{unicode-math} % this also loads fontspec
  \defaultfontfeatures{Scale=MatchLowercase}
  \defaultfontfeatures[\rmfamily]{Ligatures=TeX,Scale=1}
\fi
\usepackage{lmodern}
\ifPDFTeX\else
  % xetex/luatex font selection
\fi
% Use upquote if available, for straight quotes in verbatim environments
\IfFileExists{upquote.sty}{\usepackage{upquote}}{}
\IfFileExists{microtype.sty}{% use microtype if available
  \usepackage[]{microtype}
  \UseMicrotypeSet[protrusion]{basicmath} % disable protrusion for tt fonts
}{}
\makeatletter
\@ifundefined{KOMAClassName}{% if non-KOMA class
  \IfFileExists{parskip.sty}{%
    \usepackage{parskip}
  }{% else
    \setlength{\parindent}{0pt}
    \setlength{\parskip}{6pt plus 2pt minus 1pt}}
}{% if KOMA class
  \KOMAoptions{parskip=half}}
\makeatother
\usepackage{xcolor}
\usepackage{color}
\usepackage{fancyvrb}
\newcommand{\VerbBar}{|}
\newcommand{\VERB}{\Verb[commandchars=\\\{\}]}
\DefineVerbatimEnvironment{Highlighting}{Verbatim}{commandchars=\\\{\}}
% Add ',fontsize=\small' for more characters per line
\newenvironment{Shaded}{}{}
\newcommand{\AlertTok}[1]{\textcolor[rgb]{1.00,0.00,0.00}{\textbf{#1}}}
\newcommand{\AnnotationTok}[1]{\textcolor[rgb]{0.38,0.63,0.69}{\textbf{\textit{#1}}}}
\newcommand{\AttributeTok}[1]{\textcolor[rgb]{0.49,0.56,0.16}{#1}}
\newcommand{\BaseNTok}[1]{\textcolor[rgb]{0.25,0.63,0.44}{#1}}
\newcommand{\BuiltInTok}[1]{\textcolor[rgb]{0.00,0.50,0.00}{#1}}
\newcommand{\CharTok}[1]{\textcolor[rgb]{0.25,0.44,0.63}{#1}}
\newcommand{\CommentTok}[1]{\textcolor[rgb]{0.38,0.63,0.69}{\textit{#1}}}
\newcommand{\CommentVarTok}[1]{\textcolor[rgb]{0.38,0.63,0.69}{\textbf{\textit{#1}}}}
\newcommand{\ConstantTok}[1]{\textcolor[rgb]{0.53,0.00,0.00}{#1}}
\newcommand{\ControlFlowTok}[1]{\textcolor[rgb]{0.00,0.44,0.13}{\textbf{#1}}}
\newcommand{\DataTypeTok}[1]{\textcolor[rgb]{0.56,0.13,0.00}{#1}}
\newcommand{\DecValTok}[1]{\textcolor[rgb]{0.25,0.63,0.44}{#1}}
\newcommand{\DocumentationTok}[1]{\textcolor[rgb]{0.73,0.13,0.13}{\textit{#1}}}
\newcommand{\ErrorTok}[1]{\textcolor[rgb]{1.00,0.00,0.00}{\textbf{#1}}}
\newcommand{\ExtensionTok}[1]{#1}
\newcommand{\FloatTok}[1]{\textcolor[rgb]{0.25,0.63,0.44}{#1}}
\newcommand{\FunctionTok}[1]{\textcolor[rgb]{0.02,0.16,0.49}{#1}}
\newcommand{\ImportTok}[1]{\textcolor[rgb]{0.00,0.50,0.00}{\textbf{#1}}}
\newcommand{\InformationTok}[1]{\textcolor[rgb]{0.38,0.63,0.69}{\textbf{\textit{#1}}}}
\newcommand{\KeywordTok}[1]{\textcolor[rgb]{0.00,0.44,0.13}{\textbf{#1}}}
\newcommand{\NormalTok}[1]{#1}
\newcommand{\OperatorTok}[1]{\textcolor[rgb]{0.40,0.40,0.40}{#1}}
\newcommand{\OtherTok}[1]{\textcolor[rgb]{0.00,0.44,0.13}{#1}}
\newcommand{\PreprocessorTok}[1]{\textcolor[rgb]{0.74,0.48,0.00}{#1}}
\newcommand{\RegionMarkerTok}[1]{#1}
\newcommand{\SpecialCharTok}[1]{\textcolor[rgb]{0.25,0.44,0.63}{#1}}
\newcommand{\SpecialStringTok}[1]{\textcolor[rgb]{0.73,0.40,0.53}{#1}}
\newcommand{\StringTok}[1]{\textcolor[rgb]{0.25,0.44,0.63}{#1}}
\newcommand{\VariableTok}[1]{\textcolor[rgb]{0.10,0.09,0.49}{#1}}
\newcommand{\VerbatimStringTok}[1]{\textcolor[rgb]{0.25,0.44,0.63}{#1}}
\newcommand{\WarningTok}[1]{\textcolor[rgb]{0.38,0.63,0.69}{\textbf{\textit{#1}}}}
\setlength{\emergencystretch}{3em} % prevent overfull lines
\providecommand{\tightlist}{%
  \setlength{\itemsep}{0pt}\setlength{\parskip}{0pt}}
\setcounter{secnumdepth}{-\maxdimen} % remove section numbering
\ifLuaTeX
  \usepackage{selnolig}  % disable illegal ligatures
\fi
\usepackage{bookmark}
\IfFileExists{xurl.sty}{\usepackage{xurl}}{} % add URL line breaks if available
\urlstyle{same}
\hypersetup{
  hidelinks,
  pdfcreator={LaTeX via pandoc}}

\author{}
\date{}

\begin{document}

\section{Proposal Tugas Akhir - Versi
LaTeX}\label{proposal-tugas-akhir---versi-latex}

Dokumen ini merupakan konversi LaTeX yang dapat diedit dari
\texttt{Proposal.pdf}.

\subsection{Struktur}\label{struktur}

\begin{itemize}
\tightlist
\item
  \texttt{main.tex} - berkas utama
\item
  \texttt{frontmatter.tex} - sampul, pengesahan, orisinalitas, abstrak,
  kata pengantar, dan daftar
\item
  \texttt{chapter1.tex} - Bab I
\item
  \texttt{chapter2.tex} - Bab II
\item
  \texttt{chapter3.tex} - Bab III
\item
  \texttt{appendices.tex} - lampiran
\item
  \texttt{references.bib} - daftar pustaka format BibLaTeX/IEEE
\item
  \texttt{figures/} - logo dan gambar yang diekstrak dari PDF sumber
\end{itemize}

\subsection{Kompilasi di Overleaf}\label{kompilasi-di-overleaf}

\begin{enumerate}
\def\labelenumi{\arabic{enumi}.}
\tightlist
\item
  Unggah seluruh isi folder ini ke satu proyek Overleaf.
\item
  Pilih compiler \textbf{XeLaTeX}.
\item
  Overleaf akan menjalankan Biber secara otomatis. Jika sitasi belum
  muncul, lakukan \textbf{Recompile from scratch}.
\end{enumerate}

\subsection{Kompilasi lokal}\label{kompilasi-lokal}

\begin{Shaded}
\begin{Highlighting}[]
\ExtensionTok{xelatex}\NormalTok{ main.tex}
\ExtensionTok{biber}\NormalTok{ main}
\ExtensionTok{xelatex}\NormalTok{ main.tex}
\ExtensionTok{xelatex}\NormalTok{ main.tex}
\end{Highlighting}
\end{Shaded}

\subsection{Kolaborasi GitHub}\label{kolaborasi-github}

Folder ini belum menjadi repository Git, jadi langkah awalnya adalah
menginisialisasi repo lalu menghubungkannya ke GitHub.

\subsubsection{File yang perlu di-track}\label{file-yang-perlu-di-track}

\begin{itemize}
\tightlist
\item
  \texttt{main.tex}
\item
  \texttt{frontmatter.tex}
\item
  \texttt{chapter1.tex}
\item
  \texttt{chapter2.tex}
\item
  \texttt{chapter3.tex}
\item
  \texttt{appendices.tex}
\item
  \texttt{references.bib}
\item
  \texttt{figures/}
\item
  \texttt{README.md}
\item
  \texttt{.gitignore}
\end{itemize}

\subsubsection{File yang tidak perlu
di-track}\label{file-yang-tidak-perlu-di-track}

Hasil kompilasi LaTeX seperti \texttt{.aux}, \texttt{.log},
\texttt{.toc}, \texttt{.bbl}, \texttt{.bcf}, \texttt{.run.xml}, dan
\texttt{.pdf} sudah diabaikan melalui \texttt{.gitignore}.

\subsubsection{SOP kerja kolaboratif}\label{sop-kerja-kolaboratif}

\begin{enumerate}
\def\labelenumi{\arabic{enumi}.}
\tightlist
\item
  Pemilik repo membuat repository GitHub dan mengunggah isi folder ini.
\item
  Setiap kolaborator melakukan \texttt{git\ clone} ke laptop
  masing-masing.
\item
  Jangan bekerja langsung di branch \texttt{main}. Buat branch baru
  sesuai tugas, misalnya \texttt{revisi-bab1}, \texttt{update-abstrak},
  atau \texttt{perbaiki-sitasi}.
\item
  Lakukan perubahan pada file yang memang menjadi tanggung jawab tugas
  tersebut.
\item
  Compile lokal terlebih dahulu sebelum commit untuk memastikan tidak
  ada error.
\item
  Commit dengan pesan yang jelas, misalnya
  \texttt{revisi\ latar\ belakang\ bab\ 1} atau
  \texttt{tambah\ referensi\ bab\ 2}.
\item
  Push branch ke GitHub lalu buat Pull Request ke \texttt{main}.
\item
  Reviewer memeriksa isi perubahan, menarik branch tersebut jika perlu,
  lalu compile ulang secara lokal.
\item
  Jika hasilnya aman, Pull Request di-merge ke \texttt{main}.
\end{enumerate}

\subsubsection{Saran pembagian kerja}\label{saran-pembagian-kerja}

\begin{itemize}
\tightlist
\item
  Satu orang fokus pada satu file atau satu bab agar konflik kecil.
\item
  Jika mengubah \texttt{references.bib}, kabari anggota tim lain karena
  file ini rawan konflik.
\item
  Hindari mengedit file yang sama secara bersamaan, terutama
  \texttt{frontmatter.tex} dan \texttt{references.bib}.
\item
  Gunakan nama branch yang singkat dan jelas.
\end{itemize}

\subsubsection{Contoh alur command}\label{contoh-alur-command}

\begin{Shaded}
\begin{Highlighting}[]
\FunctionTok{git}\NormalTok{ clone }\OperatorTok{\textless{}}\NormalTok{url{-}repo}\OperatorTok{\textgreater{}}
\BuiltInTok{cd}\NormalTok{ Proposal\_LaTeX}
\FunctionTok{git}\NormalTok{ checkout }\AttributeTok{{-}b}\NormalTok{ revisi{-}bab2}
\end{Highlighting}
\end{Shaded}

Setelah selesai mengedit dan berhasil compile:

\begin{Shaded}
\begin{Highlighting}[]
\FunctionTok{git}\NormalTok{ add .}
\FunctionTok{git}\NormalTok{ commit }\AttributeTok{{-}m} \StringTok{"revisi tinjauan pustaka bab 2"}
\FunctionTok{git}\NormalTok{ push origin revisi{-}bab2}
\end{Highlighting}
\end{Shaded}


\end{document}
